\ifdefined\gkver\else\def\gkver{student}\fi
\documentclass[\gkver]{gaokaozhenti}
\begin{document}

\chapter{2023年广东}

\begin{enumerate}
\item
%% number: 1
%% paperName: 2023年广东高考第 1 题 4 分
%% typeId: 1
%% type: 单选题
%% chapter: 原子和原子核
%% point: 人工核反应
%% method: 无
%% score: 4
%% degree: 950
%% duplicateId: 0
%% body:
理论认为，大质量恒星塌缩成黑洞的过程，受核反应 $\ce{^{12}_{6}C + Y -> ^{16}_{8}O}$ 的影响。下列说法正确的是\xzanswer{D}

\fourchoices[answer=D]
{Y 是 $\alpha$ 粒子，$\beta$ 射线穿透能力比 $\gamma$ 射线强}
{Y 是 $\beta$ 粒子，$\beta$ 射线电离能力比 $\gamma$ 射线强}
{Y 是 $\alpha$ 粒子，$\alpha$ 射线穿透能力比 $\gamma$ 射线强}
{Y 是 $\alpha$ 粒子，$\alpha$ 射线电离能力比 $\gamma$ 射线强}

%% answer: D
%% memo:
\memoanswer{
根据受核反应满足质量数和电荷数守恒可知，Y 是 $\alpha$ 粒子（$\ce{^{4}_{2}He}$），三种射线的穿透能力，$\gamma$ 射线最强，$\alpha$ 射线最弱；三种射线的电离能力，$\alpha$ 射线最强，$\gamma$ 射线最弱。

故选 D。
}

\item
%% number: 2
%% paperName: 2023年广东高考第 2 题 4 分
%% typeId: 1
%% type: 单选题
%% chapter: 力物体的平衡
%% point: 多力平衡
%% method: 无
%% score: 4
%% degree: 650
%% duplicateId: 0
%% body:
如图所示，可视为质点的机器人通过磁铁吸附在船舷外壁面检测船体。壁面可视为斜面，与竖直方向夹角为 $\theta$。船和机器人保持静止时，机器人仅受重力 $G$、支持力 $F_{N}$、摩擦力 $F_{f}$ 和磁力 $F$ 的作用，磁力垂直壁面。下列关系式正确的是\xzanswer{C}

\onepicture[w=3.20cm]{figs/02.png}

\fourchoices[answer=C]
{$F_{f} = G$}
{$F = F_{N}$}
{$F_{f} = G\cos\theta$}
{$F = G\sin\theta$}

%% answer: C
%% memo:
\memoanswer{
如图所示，将重力垂直于斜面方向和沿斜面方向分解。

\onepicture[w=3.20cm]{figs/02_an.png}

AC．沿斜面方向，由平衡条件得 $F_{f} = G\cos\theta$

故 A 错误，C 正确；

BD．垂直斜面方向，由平衡条件得

$F = G\sin\theta + F_{N}$

故 BD 错误。

故选 C。
}

\item
%% number: 3
%% paperName: 2023年广东高考第 3 题 4 分
%% typeId: 1
%% type: 单选题
%% chapter: 直线运动
%% point: 竖直上下抛运动
%% method: 无
%% score: 4
%% degree: 750
%% duplicateId: 0
%% body:
铯原子喷泉钟是定标“秒”的装置。在喷泉钟的真空系统中，可视为质点的铯原子团在激光的推动下，获得一定的初速度。随后激光关闭，铯原子团仅在重力的作用下做竖直上抛运动，到达最高点后再做一段自由落体运动。取竖直向上为正方向。下列可能表示激光关闭后铯原子团速度 $v$ 或加速度 $a$ 随时间 $t$ 变化的图像是\xzanswer{D}

\fourchoices[answer=D, ispicture=true, h=2.20cm]
{figs/03a.png}
{figs/03b.png}
{figs/03c.png}
{figs/03d.png}

%% answer: D
%% memo:
\memoanswer{
AB．铯原子团仅在重力的作用，加速度 $g$ 竖直向下，大小恒定，在 $v - t$ 图像中，斜率为加速度，故斜率不变，所以图像应该是一条倾斜的直线，故选项 AB 错误；

CD．因为加速度恒定，且方向竖直向下，故为负值，故选项 C 错误，选项 D 正确。

故选 D。
}

\item
%% number: 4
%% paperName: 2023年广东高考第 4 题 4 分
%% typeId: 1
%% type: 单选题
%% chapter: 机械波
%% point: 波长频率波速
%% method: 无
%% score: 4
%% degree: 650
%% duplicateId: 0
%% body:
渔船常用回声探测器发射的声波探测水下鱼群与障碍物，声波在水中传播速度为 $1500 \Ums$，若探测器发出频率为 $1.5\times10^{6} \UHz$ 的声波，下列说法正确的是\xzanswer{B}

\fourchoices[answer=B]
{两列声波相遇时一定会发生干涉}
{声波由水中传播到空气中，波长会改变}
{该声波遇到尺寸约为 $1 \Um$ 的被探测物时会发生明显衍射}
{探测器接收到的回声频率与被探测物相对探测器运动的速度无关}

%% answer: B
%% memo:
\memoanswer{
AD．根据多普勒效应可知，探测器接收到的回声频率与被探测物相对探测器运动的速度有关，而两列声波发生干涉的条件是频率相等，所以两列声波相遇时不一定发生干涉，故 AD 错误；

B．声波由水中传播到空气中时，声波的波速发生变化，所以波长会发生改变，故 B 正确；

C．根据波长的计算公式可得

$\lambda = \frac{v}{f} = \frac{1500}{1.5\times10^{6}} \Um = 1\times10^{-3} \Um$

当遇到尺寸约 $1 \Um$ 的被探测物时不会发生明显衍射，故 C 错误；

故选 B。
}

\item
%% number: 5
%% paperName: 2023年广东高考第 5 题 4 分
%% typeId: 1
%% type: 单选题
%% chapter: 磁场
%% point: 质谱仪回旋加速器
%% method: 无
%% score: 4
%% degree: 650
%% duplicateId: 0
%% body:
某小型医用回旋加速器，最大回旋半径为 $0.5 \Um$，磁感应强度大小为 $1.12 \UT$，质子加速后获得的最大动能为 $1.5\times10^{7} \UeV$。根据给出的数据，可计算质子经该回旋加速器加速后的最大速率约为（忽略相对论效应，$1 \UeV = 1.6\times10^{-19} \UJ$）\xzanswer{C}

\fourchoices[answer=C]
{$3.6\times10^{6} \Ums$}
{$1.2\times10^{7} \Ums$}
{$5.4\times10^{7} \Ums$}
{$2.4\times10^{8} \Ums$}

%% answer: C
%% memo:
\memoanswer{
洛伦兹力提供向心力有 $qvB = m\frac{v^{2}}{R}$

质子加速后获得的最大动能为 $E_{k} = \frac{1}{2}mv^{2}$

解得最大速率约为 $v = 5.4\times10^{7} \Ums$

故选 C。
}

\item
%% number: 6
%% paperName: 2023年广东高考第 6 题 4 分
%% typeId: 1
%% type: 单选题
%% chapter: 交流电
%% point: 变压器
%% method: 无
%% score: 4
%% degree: 650
%% duplicateId: 0
%% body:
用一台理想变压器对电动汽车充电，该变压器原、副线圈的匝数比为 $1:2$，输出功率为 $8.8 \UkW$，原线圈的输入电压 $u = 220\sqrt{2}\sin 100\pi t \UV$。关于副线圈输出电流的有效值和频率正确的是\xzanswer{A}

\fourchoices[answer=A]
{$20 \UA$，$50 \UHz$}
{$20\sqrt{2} \UA$，$50 \UHz$}
{$20 \UA$，$100 \UHz$}
{$20\sqrt{2} \UA$，$100 \UHz$}

%% answer: A
%% memo:
\memoanswer{
由题可知原线圈输入电压的有效值为

$U_{1} = \frac{220\sqrt{2}}{\sqrt{2}} \UV = 220 \UV$

原线圈电流为

$I_{1} = \frac{P}{U_{1}} = \frac{8.8\times10^{3}}{220} \UA = 40 \UA$

副线圈输出电流的有效值为

$I_{2} = \frac{n_{1}}{n_{2}}I_{1} = 20 \UA$

变压器无法改变电流的频率，故

$f = \frac{\omega}{2\pi} = \frac{100\pi}{2\pi} \UHz = 50 \UHz$

故选 A。
}

\item
%% number: 7
%% paperName: 2023年广东高考第 7 题 6 分
%% typeId: 1
%% type: 单选题
%% chapter: 曲线运动
%% point: 天体运动
%% method: 无
%% score: 6
%% degree: 450
%% duplicateId: 0
%% body:
如图 \ref{2023广东07a} 所示，太阳系外的一颗行星 P 绕恒星 Q 做匀速圆周运动。由于 P 的遮挡，探测器探测到 Q 的亮度随时间做如图 \ref{2023广东07b} 所示的周期性变化，该周期与 P 的公转周期相同。已知 Q 的质量为 $M$，引力常量为 $G$。关于 P 的公转，下列说法正确的是\xzanswer{B}

\twopicture[labelA=2023广东07a, labelB=2023广东07b, wA=5.00cm, wB=5.50cm]
{figs/07a.png}{figs/07b.png}

\fourchoices[answer=B]
{周期为 $2t_{1} - t_{0}$}
{半径为 $\sqrt[3]{\frac{GM(t_{1} - t_{0})^{2}}{4\pi^{2}}}$}
{角速度的大小为 $\frac{\pi}{t_{1} - t_{0}}$}
{加速度的大小为 $\sqrt[3]{\frac{2\pi GM}{t_{1} - t_{0}}}$}

%% answer: B
%% memo:
\memoanswer{
A．由图（b）可知探测器探测到 Q 的亮度随时间变化的周期为 $T = t_{1} - t_{0}$，则 P 的公转周期为 $t_{1} - t_{0}$，故 A 错误；

B．P 绕恒星 Q 做匀速圆周运动，由万有引力提供向心力可得

$\frac{GMm}{r^{2}} = m\frac{4\pi^{2}}{T^{2}}r$

解得半径为 $r = \sqrt[3]{\frac{GMT^{2}}{4\pi^{2}}} = \sqrt[3]{\frac{GM(t_{1} - t_{0})^{2}}{4\pi^{2}}}$

故 B 正确；

C．P 的角速度为 $\omega = \frac{2\pi}{T} = \frac{2\pi}{t_{1} - t_{0}}$

故 C 错误；

D．P 的加速度大小为

$a = \omega^{2}r = \left(\frac{2\pi}{t_{1} - t_{0}}\right)^{2} \cdot \sqrt[3]{\frac{GM(t_{1} - t_{0})^{2}}{4\pi^{2}}} = \frac{2\pi}{t_{1} - t_{0}} \cdot \sqrt[3]{\frac{2\pi GM}{t_{1} - t_{0}}}$

故 D 错误。

故选 B。
}

\item
%% number: 8
%% paperName: 2023年广东高考第 8 题 6 分
%% typeId: 2
%% type: 多选题
%% chapter: 曲线运动
%% point: 竖直平面圆周运动
%% method: 无
%% score: 6
%% degree: 650
%% duplicateId: 0
%% body:
人们用滑道从高处向低处运送货物。如图所示，可看作质点的货物从 $\frac{1}{4}$ 圆弧滑道顶端 P 点静止释放，沿滑道运动到圆弧末端 Q 点时速度大小为 $6 \Ums$。已知货物质量为 $20 \Ukg$，滑道高度 $h$ 为 $4 \Um$，且过 Q 点的切线水平，重力加速度取 $10 \Umsq$。关于货物从 P 点运动到 Q 点的过程，下列说法正确的有\xzanswer{BCD}

\onepicture[w=4.20cm]{figs/08.png}

\fourchoices[answer=BCD]
{重力做的功为 $360 \UJ$}
{克服阻力做的功为 $440 \UJ$}
{经过 Q 点时向心加速度大小为 $9 \Umsq$}
{经过 Q 点时对轨道的压力大小为 $380 \UN$}

%% answer: BCD
%% memo:
\memoanswer{
A．重力做的功为 $W_{G} = mgh = 800 \UJ$。A 错误；

B．下滑过程据动能定理可得 $W_{G} - W_{f} = \frac{1}{2}mv_{Q}^{2}$，代入数据解得，克服阻力做的功为 $W_{f} = 440 \UJ$。B 正确；

C．经过 Q 点时向心加速度大小为 $a = \frac{v_{Q}^{2}}{h} = 9 \Umsq$。C 正确；

D．经过 Q 点时，据牛顿第二定律可得 $F - mg = ma$

解得货物受到的支持力大小为 $F = 380 \UN$

据牛顿第三定律可知，货物对轨道的压力大小为 $380 \UN$，D 正确。

故选 BCD。
}

\item
%% number: 9
%% paperName: 2023年广东高考第 9 题 6 分
%% typeId: 2
%% type: 多选题
%% chapter: 电场
%% point: 电势能电势差电场力做功
%% method: 无
%% score: 6
%% degree: 550
%% duplicateId: 0
%% body:
电子墨水是一种无光源显示技术，它利用电场调控带电颜料微粒的分布，使之在自然光的照射下呈现出不同颜色。透明面板下有一层胶囊，其中每个胶囊都是一个像素。如图所示，胶囊中有带正电的白色微粒和带负电的黑色微粒。当胶囊下方的电极极性由负变正时，微粒在胶囊内迁移（每个微粒电量保持不变），像素由黑色变成白色。下列说法正确的有\xzanswer{AC}

\twopicture[showlabel=false, wA=5.80cm, wB=3.05cm]
{figs/09a.png}{figs/09b.png}

\fourchoices[answer=AC]
{像素呈黑色时，黑色微粒所在区域的电势高于白色微粒所在区域的电势}
{像素呈白色时，黑色微粒所在区域的电势低于白色微粒所在区域的电势}
{像素由黑变白的过程中，电场力对白色微粒做正功}
{像素由白变黑的过程中，电场力对黑色微粒做负功}

%% answer: AC
%% memo:
\memoanswer{
A．像素呈黑色时，当胶囊下方的电极带负电，像素胶囊里电场线方向向下，所以黑色微粒所在的区域的电势高于白色微粒所在区域的电势，故 A 正确；

B．像素呈白色时，当胶囊下方的电极带正电，像素胶囊里电场线方向向上，所以黑色微粒所在的区域的电势高于白色微粒所在区域的电势，故 B 错误；

C．像素由黑变白的过程中，白色微粒受到的电场力向上，位移向上，电场力对白色微粒做正功，故 C 正确；

D．像素由白变黑的过程中，黑色微粒受到的电场力向下，位移向下，电场力对黑色微粒做正功，故 D 错误。

故选 AC。
}

\item
%% number: 10
%% paperName: 2023年广东高考第 10 题 6 分
%% typeId: 2
%% type: 多选题
%% chapter: 运动和力
%% point: 动量守恒定律
%% method: 无
%% score: 6
%% degree: 450
%% duplicateId: 0
%% body:
某同学受电动窗帘的启发，设计了如图所示的简化模型。多个质量均为 $1 \Ukg$ 的滑块可在水平滑轨上滑动，忽略阻力。开窗帘过程中，电机对滑块 1 施加一个水平向右的恒力 $F$，推动滑块 1 以 $0.40 \Ums$ 的速度与静止的滑块 2 碰撞，碰撞时间为 $0.04 \Us$，碰撞结束后瞬间两滑块的共同速度为 $0.22 \Ums$。关于两滑块的碰撞过程，下列说法正确的有\xzanswer{BD}

\onepicture[w=5.50cm]{figs/10.png}

\fourchoices[answer=BD]
{该过程动量守恒}
{滑块 1 受到合外力的冲量大小为 $0.18 \UNs$}
{滑块 2 受到合外力的冲量大小为 $0.40\UN\cdot a$}
{滑块 2 受到滑块 1 的平均作用力大小为 $5.5 \UN$}

%% answer: BD
%% memo:
\memoanswer{
A．取向右为正方向，滑块 1 和滑块 2 组成的系统的初动量为

$p_{1} = mv_{1} = 1\times0.40 \Ukgms = 0.40 \Ukgms$

碰撞后的动量为

$p_{2} = 2mv_{2} = 2\times1\times0.22 \Ukgms = 0.44 \Ukgms$

则滑块的碰撞过程动量不守恒，故 A 错误；

B．对滑块 1，取向右为正方向，则有

$I_{1} = mv_{2} - mv_{1} = (1\times0.22 - 1\times0.40) \Ukgms = -0.18 \Ukgms$

负号表示方向水平向左，故 B 正确；

C．对滑块 2，取向右为正方向，则有

$I_{2} = mv_{2} = 1\times0.22 \Ukgms = 0.22 \Ukgms$

故 C 错误；

D．对滑块 2 根据动量定理有 $F\Delta t = I_{2}$

解得 $F = 5.5 \UN$

则滑块 2 受到滑块 1 的平均作用力大小为 $5.5 \UN$，故 D 正确。

故选 BD。
}

\item
%% number: 11
%% paperName: 2023年广东高考第 11 题 7 分
%% typeId: 4
%% type: 实验题
%% chapter: 光学
%% point: 测定玻璃折射率
%% method: 无
%% score: 7
%% degree: 550
%% duplicateId: 0
%% body:
某同学用激光笔和透明长方体玻璃砖测量玻璃的折射率，实验过程如下：

\begin{enumerate}
	\item
	将玻璃砖平放在水平桌面上的白纸上，用大头针在白纸上标记玻璃砖的边界。
	\item
	\begin{steps}
		\item
		激光笔发出的激光从玻璃砖上的 M 点水平入射，到达 ef 面上的 O 点后反射到 N 点射出，用大头针在白纸上标记 O 点、M 点和激光笔出光孔 Q 的位置；
		\item
		移走玻璃砖，在白纸上描绘玻璃砖的边界和激光的光路，作 QM 连线的延长线与 ef 面的边界交于 P 点，如图 \ref{2023广东11a} 所示；
		\item
		用刻度尺测量 PM 和 OM 的长度 $d_{1}$ 和 $d_{2}$，PM 的示数如图 \ref{2023广东11b} 所示，$d_{1}$ 为\tkanswer[1.5cm]{2.25} $\Ucm$。测得 $d_{2}$ 为 $3.40 \Ucm$。
	\end{steps}
	\item
	利用所测量的物理量，写出玻璃砖折射率的表达式 $n$ = \tkanswer[1.8cm]{$\frac{d_{2}}{d_{1}}$}；由测得的数据可得折射率 $n$ 为\tkanswer[1.2cm]{1.51}（结果保留 3 位有效数字）。
	\item
	相对误差的计算式为 $\delta = \frac{\text{测量值} - \text{真实值}}{\text{真实值}} \times 100\%$。为了减小 $d_{1}$、$d_{2}$ 测量的相对误差，实验中激光在 M 点入射时应尽量使入射角\tkanswer[1.5cm]{稍小一些}。
\end{enumerate}

\twopicture[labelA=2023广东11a, labelB=2023广东11b, wA=4.20cm, wB=6.50cm]
{figs/11a.png}{figs/11b.png}

%% answer:
\jdanswer{
\begin{enumerate}
	\item
	—
	\item
	$2.25$
	\item
	$\frac{d_{2}}{d_{1}}$，$1.51$
	\item
	稍小一些
\end{enumerate}
}
%% memo:
\memoanswer{
（2）③刻度尺的最小分度为 $0.1 \Ucm$，由图可知，$d_{1}$ 为 $2.25 \Ucm$；

（3）玻璃砖折射率的表达式

$n = \frac{\sin i}{\sin r} = \frac{\frac{fM}{MP}}{\frac{fM}{OM}} = \frac{OM}{PM} = \frac{d_{2}}{d_{1}}$

代入数据可知 $n = \frac{3.40}{2.25} = 1.51$

（4）为了减小 $d_{1}$、$d_{2}$ 测量的相对误差，实验中 $d_{1}$、$d_{2}$ 要尽量稍大一些，即激光在 M 点入射时应尽量使入射角稍小一些。
}

\item
%% number: 12
%% paperName: 2023年广东高考第 12 题 10 分
%% typeId: 4
%% type: 实验题
%% chapter: 电路
%% point: 测量电阻率
%% method: 无
%% score: 10
%% degree: 450
%% duplicateId: 0
%% body:
某兴趣小组设计了测量盐水电导率的实验。所用器材有：电源 $E$（电动势恒定，内阻可忽略）；毫安表（量程 $15 \UmA$，内阻可忽略）；电阻 $R_{1}$（阻值 $500 \UO$）、$R_{2}$（阻值 $500 \UO$）、$R_{3}$（阻值 $600 \UO$）和 $R_{4}$（阻值 $200 \UO$）；开关 $\mathrm{S}_{1}$ 和 $\mathrm{S}_{2}$；装有耐腐蚀电极板和温度计的有机玻璃样品池；导线若干。请完成下列实验操作和计算。

\begin{enumerate}
	\item
	图 \ref{2023广东12a} 为实验原理图。在图 \ref{2023广东12b} 的实物图中，已正确连接了部分电路，只有 $R_{4}$ 一端的导线还末连接，该导线应接到 $R_{3}$ 的\tkanswer[1.0cm]{右}（填“左”或“右”）端接线柱。
	\item
	\begin{steps}
		\item
		测量并记录样品池内壁的长宽高。在样品池中注满待测盐水。
		\item
		闭合开关 $\mathrm{S}_{1}$，\tkanswer[1.0cm]{断开}开关 $\mathrm{S}_{2}$，毫安表的示数为 $10.0 \UmA$，记录此时毫安表的示数；计算得到流过样品池的电流 $I_{1}$ 为\tkanswer[1.2cm]{40.0} $\UmA$。
		\item
		\tkanswer[1.0cm]{闭合}开关 $\mathrm{S}_{2}$，毫安表的示数为 $15.0 \UmA$，记录此时毫安表的示数；计算得到流过样品池的电流 $I_{2}$ 为\tkanswer[1.2cm]{60.0} $\UmA$。
		\item
		断开开关 $\mathrm{S}_{1}$，测量并记录盐水的温度。
	\end{steps}
	\item
	根据上述数据，计算得到样品池两电极板间待测盐水的电阻为\tkanswer[1.2cm]{100} $\UO$，进而可求得该温度时待测盐水的电导率。
\end{enumerate}

\twopicture[labelA=2023广东12a, labelB=2023广东12b, wA=5.00cm, wB=7.00cm]
{figs/12a.png}{figs/12b.png}

%% answer:
\jdanswer{
\begin{enumerate}
	\item
	右
	\item
	②断开，$40.0$；③闭合，$60.0$
	\item
	$100$
\end{enumerate}
}
%% memo:
\memoanswer{
（1）根据图（a）为电路可知，$R_{4}$ 一端的导线应接到 $R_{3}$ 的右端接线柱；

（2）②闭合开关 $\mathrm{S}_{1}$，断开开关 $\mathrm{S}_{2}$，毫安表的示数为 $10.0 \UmA$，则通过电阻 $R_{4}$ 的电流为 $I_{4} = \frac{I_{3}R_{3}}{R_{4}}$

根据电路构造可知，流过样品池的电流为

$I_{1} = I_{3} + I_{4} = 40.0 \UmA$

③闭合开关 $\mathrm{S}_{2}$，毫安表的示数为 $15.0 \UmA$，则流过 $R_{4}$ 的电流为

$I_{4}' = \frac{I_{3}'R_{3}}{R_{4}}$

流过样品池的电流 $I_{2}$ 为

$I_{2} = I_{3}' + I_{4}' = 60.0 \UmA$

（3）设待测盐水的电阻为 $R_{0}$，根据闭合电路欧姆定律，开关 $\mathrm{S}_{2}$ 断开时

$E = I_{1}\left(R_{0} + R_{2} + \frac{R_{3}R_{4}}{R_{3} + R_{4}}\right)$

开关 $\mathrm{S}_{2}$ 闭合时

$E = I_{2}\left(R_{0} + \frac{R_{1}R_{2}}{R_{1} + R_{2}} + \frac{R_{3}R_{4}}{R_{3} + R_{4}}\right)$

代入数据解得 $R_{0} = 100 \UO$
}

\item
%% number: 13
%% paperName: 2023年广东高考第 13 题 9 分
%% typeId: 6
%% type: 计算题
%% chapter: 气体性质
%% point: 气体多过程
%% method: 无
%% score: 9
%% degree: 550
%% duplicateId: 0
%% body:
在驻波声场作用下，水中小气泡周围液体的压强会发生周期性变化，使小气泡周期性膨胀和收缩，气泡内气体可视为质量不变的理想气体，其膨胀和收缩过程可简化为如图所示的 $p - V$ 图像，气泡内气体先从压强为 $p_{0}$、体积为 $V_{0}$、温度为 $T_{0}$ 的状态 A 等温膨胀到体积为 $5V_{0}$、压强为 $p_{B}$ 的状态 B，然后从状态 B 绝热收缩到体积为 $V_{0}$、压强为 $1.9p_{0}$、温度为 $T_{C}$ 的状态 C，B 到 C 过程中外界对气体做功为 $W$。已知 $p_{0}$、$V_{0}$、$T_{0}$ 和 $W$。求：

\begin{enumerate}
	\item
	$p_{B}$ 的表达式；
	\item
	$T_{C}$ 的表达式；
	\item
	B 到 C 过程，气泡内气体的内能变化了多少？
\end{enumerate}

\onepicture[w=4.00cm, align=right]{figs/13.png}

%% answer:
\jdanswer{
\begin{enumerate}
	\item
	$p_{B} = \frac{1}{5}p_{0}$
	\item
	$T_{C} = 1.9T_{0}$
	\item
	$\Delta U = W$
\end{enumerate}
}
%% memo:
\memoanswer{
（1）由题可知，根据玻意耳定律可得 $p_{A}V_{A} = p_{B}V_{B}$

解得 $p_{B} = \frac{1}{5}p_{0}$

（2）根据理想气体状态方程可知 $\frac{p_{B}V_{B}}{T_{B}} = \frac{p_{C}V_{C}}{T_{C}}$

解得 $T_{C} = 1.9T_{0}$

（3）根据热力学第一定律可知 $\Delta U = W + Q$，其中 $Q = 0$，故气体内能增加 $\Delta U = W$。
}

\item
%% number: 14
%% paperName: 2023年广东高考第 14 题 13 分
%% typeId: 6
%% type: 计算题
%% chapter: 电磁感应
%% point: 线圈过有界磁场
%% method: 无
%% score: 13
%% degree: 450
%% duplicateId: 0
%% body:
光滑绝缘的水平面上有垂直平面的匀强磁场，磁场被分成区域 Ⅰ 和 Ⅱ，宽度均为 $h$，其俯视图如图 \ref{2023广东14a} 所示，两磁场磁感应强度随时间 $t$ 的变化如图 \ref{2023广东14b} 所示，$0 \sim \tau$ 时间内，两区域磁场恒定，方向相反，磁感应强度大小分别为 $2B_{0}$ 和 $B_{0}$，一电阻为 $R$、边长为 $h$ 的刚性正方形金属框 abcd，平放在水平面上，ab、cd 边与磁场边界平行。$t = 0$ 时，线框 ab 边刚好跨过区域 Ⅰ 的左边界以速度 $v$ 向右运动。在 $\tau$ 时刻，ab 边运动到距区域 Ⅰ 的左边界 $\frac{h}{2}$ 处，线框的速度近似为零，此时线框被固定，如图 \ref{2023广东14a} 中的虚线框所示。随后在 $\tau \sim 2\tau$ 时间内，Ⅰ 区磁感应强度线性减小到 0，Ⅱ 区磁场保持不变；$2\tau \sim 3\tau$ 时间内，Ⅱ 区磁感应强度也线性减小到 0。求：

\begin{enumerate}
	\item
	$t = 0$ 时线框所受的安培力 $F$；
	\item
	$t = 1.2\tau$ 时穿过线框的磁通量 $\Phi$；
	\item
	$2\tau \sim 3\tau$ 时间内，线框中产生的热量 $Q$。
\end{enumerate}

\twopicture[labelA=2023广东14a, labelB=2023广东14b, wA=3.50cm, wB=6.00cm]
{figs/14a.png}{figs/14b.png}

%% answer:
\jdanswer{
\begin{enumerate}
	\item
	$F = \frac{9B_{0}^{2}h^{2}v}{R}$，方向水平向左
	\item
	$\Phi = \frac{1}{2}B_{0}h^{2}$
	\item
	$Q = \frac{B_{0}^{2}h^{4}}{4\tau R}$
\end{enumerate}
}
%% memo:
\memoanswer{
（1）由图可知 $t = 0$ 时线框切割磁感线的感应电动势为

$E = 2B_{0}hv + B_{0}hv = 3B_{0}hv$

则感应电流大小为 $I = \frac{E}{R} = \frac{3B_{0}hv}{R}$

所受的安培力为 $F = 2B_{0}\frac{3B_{0}hv}{R}h + B_{0}\frac{3B_{0}hv}{R}h = \frac{9B_{0}^{2}h^{2}v}{R}$

方向水平向左；

（2）在 $\tau$ 时刻，ab 边运动到距区域 Ⅰ 的左边界 $\frac{h}{2}$ 处，线框的速度近似为零，此时线框被固定，则 $t = 1.2\tau$ 时穿过线框的磁通量为

$\Phi = 2B_{0}h \cdot \frac{1}{2}h - B_{0}h \cdot \frac{1}{2}h = \frac{1}{2}B_{0}h^{2}$

方向垂直纸面向里；

（3）$2\tau \sim 3\tau$ 时间内，Ⅱ 区磁感应强度也线性减小到 0，则有

$E' = \frac{\Delta\Phi}{\Delta t} = \frac{B_{0}\frac{1}{2}h^{2}}{\tau} = \frac{B_{0}h^{2}}{2\tau}$

感应电流大小为

$I' = \frac{E'}{R} = \frac{B_{0}h^{2}}{2\tau R}$

则 $2\tau \sim 3\tau$ 时间内，线框中产生的热量为

$Q = I'^{2}Rt = \frac{B_{0}^{2}h^{4}}{4\tau R}$
}

\item
%% number: 15
%% paperName: 2023年广东高考第 15 题 15 分
%% typeId: 6
%% type: 计算题
%% chapter: 运动和力
%% point: 动量和能量
%% method: 无
%% score: 15
%% degree: 350
%% duplicateId: 0
%% body:
如图为某药品自动传送系统的示意图。该系统由水平传送带、竖直螺旋滑槽和与滑槽平滑连接的平台组成，滑槽高为 $3L$，平台高为 $L$。药品盒 A、B 依次被轻放在以速度 $v_{0}$ 匀速运动的传送带上，在与传送带达到共速后，从 M 点进入滑槽，A 刚好滑到平台最右端 N 点停下，随后滑下的 B 以 $2v_{0}$ 的速度与 A 发生正碰，碰撞时间极短，碰撞后 A、B 恰好落在桌面上圆盘内直径的两端。已知 A、B 的质量分别为 $m$ 和 $2m$，碰撞过程中损失的能量为碰撞前瞬间总动能的 $\frac{1}{4}$。A 与传送带间的动摩擦因数为 $\mu$，重力加速度为 $g$，AB 在滑至 N 点之前不发生碰撞，忽略空气阻力和圆盘的高度，将药品盒视为质点。求：

\begin{enumerate}
	\item
	A 在传送带上由静止加速到与传送带共速所用的时间 $t$；
	\item
	B 从 M 点滑至 N 点的过程中克服阻力做的功 $W$；
	\item
	圆盘的圆心到平台右端 N 点的水平距离 $s$。
\end{enumerate}

\onepicture[w=4.50cm, align=right]{figs/15.png}

%% answer:
\jdanswer{
\begin{enumerate}
	\item
	$t = \frac{v_{0}}{\mu g}$
	\item
	$W = 6mgL - 3mv_{0}^{2}$
	\item
	$s = \frac{3v_{0}}{2}\sqrt{\frac{2L}{g}}$
\end{enumerate}
}
%% memo:
\memoanswer{
（1）A 在传送带上运动时的加速度 $a = \mu g$

由静止加速到与传送带共速所用的时间 $t = \frac{v_{0}}{a} = \frac{v_{0}}{\mu g}$

（2）B 从 M 点滑至 N 点的过程中克服阻力做的功

$W = \frac{1}{2} \cdot 2mv_{0}^{2} + 2\Umg \cdot 3L - \frac{1}{2} \cdot 2m(2v_{0})^{2} = 6mgL - 3mv_{0}^{2}$

（3）AB 碰撞过程由动量守恒定律和能量关系可知

$2m \cdot 2v_{0} = mv_{1} + 2mv_{2}$

$\frac{1}{2} \cdot 2m \cdot (2v_{0})^{2} - \left(\frac{1}{2}mv_{1}^{2} + \frac{1}{2} \cdot 2mv_{2}^{2}\right) = \frac{1}{4}\left[\frac{1}{2} \cdot 2m \cdot (2v_{0})^{2}\right]$

解得 $v_{1} = 2v_{0}$，$v_{2} = v_{0}$（另一组 $v_{1} = \frac{2}{3}v_{0}$，$v_{2} = \frac{5}{3}v_{0}$ 舍掉）

两物体平抛运动的时间 $t_{1} = \sqrt{\frac{2L}{g}}$

则 $s - r = v_{2}t_{1}$，$s + r = v_{1}t_{1}$

解得 $s = \frac{3v_{0}}{2}\sqrt{\frac{2L}{g}}$
}

\end{enumerate}

\end{document}
